\chapter{โปรโตคอลการตรวจวิเคราะห์และสกัดดีเอ็นเอจากดิน}
\label{app:protocols}

\section{โปรโตคอลการสกัดจีโนมิกดีเอ็นเอรวมจากดิน (Soil DNA Extraction)}
การสกัด DNA จากดินมีความท้าทายเนื่องจากมีสารฮิวมัส (Humic substances) ซึ่งเป็นตัวยับยั้งเอนไซม์ Taq polymerase ในปฏิกิริยา PCR ขั้นตอนมาตรฐานประกอบด้วย:
\begin{enumerate}
    \item \textbf{การสลายเซลล์เชิงกลและเคมี (Cell Lysis)} ชั่งตัวอย่างดิน 0.25--0.50 กรัม ใส่ลงในหลอดสลายเซลล์ที่มีลูกปัดซิลิกา (Bead beating tube) เติม Lysis Buffer ที่มีสารซักฟอก SDS และสารยับยั้งเอนไซม์ EDTA นำไปเขย่าด้วยเครื่องบีดบีตเตอร์ (Bead Beater) ที่ความเร็ว 6.0 m/s เป็นเวลา 40 วินาที เพื่อสลายผนังเซลล์ของแบคทีเรียและสปอร์รา
    \item \textbf{การกำจัดสารฮิวมัสและโปรตีน (Inhibitor Removal)} เติมสารละลายตกตะกอนโปรตีนและฮิวมิกแอซิด (Potassium acetate หรือ Ammonium acetate) บ่มที่ 4 องศาเซลเซียส เป็นเวลา 10 นาที จากนั้นปั่นเหวี่ยงที่ 12,000 rpm เป็นเวลา 5 นาที เพื่อแยกตะกอนสิ่งเจือปนออก
    \item \textbf{การจับดีเอ็นเอกับเมมเบรนซิลิกา (DNA Binding)} ดูดส่วนใสผสมกับสารละลาย Binding Buffer ที่มีเกลือคาโอโทรปิก (Chaotropic salts) ส่งผ่านคอลัมน์ซิลิกา (Silica-spin column) ดีเอ็นเอจะจับกับแผ่นกรองซิลิกาอย่างจำเพาะ
    \item \textbf{การล้างสิ่งเจือปน (Washing)} ล้างคอลัมน์ด้วยสารละลาย 70\% เอทานอล 2 ครั้ง เพื่อชะล้างเกลือและสารตกค้าง
    \item \textbf{การชะล้างดีเอ็นเอบริสุทธิ์ (Elution)} เติมบัฟเฟอร์ TE หรือน้ำปลอดเชื้อเกรดโมเลกุลปริมาตร 50 ไมโครลิตร บ่มที่อุณหภูมิห้อง 2 นาที ปั่นเหวี่ยงเพื่อชะล้างดีเอ็นเอบริสุทธิ์ นำไปวัดค่าความเข้มข้นและความบริสุทธิ์ด้วยเครื่อง Spectrophotometer ($A_{260}/A_{280} \approx 1.8$ และ $A_{260}/A_{230} > 1.7$)
\end{enumerate}

\section{โปรโตคอลการตรวจนับจุลินทรีย์ดิน (Serial Dilution \& Plate Count)}
\begin{enumerate}
    \item ชั่งดิน 10 กรัม ผสมในน้ำเกลือปลอดเชื้อ (0.85\% NaCl) ปริมาตร 90 mL เขย่าเป็นเวลา 15 นาที
    \item ดูดสารแขวนลอย 1 mL ถ่ายลงในหลอดที่ 2 ที่มีน้ำปลอดเชื้อ 9 mL (เจือจาง $10^{-2}$) ทำซ้ำจนถึง $10^{-6}$
    \item ดูดตัวอย่าง 0.1 mL จากระดับ $10^{-4}, 10^{-5}, 10^{-6}$ ลงบนจานอาหาร NA (สำหรับแบคทีเรีย) และจาก $10^{-2}, 10^{-3}$ ลงบนจานอาหาร PDA ผสมยาปฏิชีวนะ (สำหรับเชื้อรา)
    \item ใช้แท่งแก้วรูปสามเหลี่ยมเกลี่ยให้ทั่วจานอาหาร บ่มที่อุณหภูมิ 30 องศาเซลเซียส ตรวจนับโคโลนีช่วง 30--300 โคโลนี
\end{enumerate}

\clearpage